% ============================================================
% chapters/ch0_introduction/introduction.tex
% ============================================================
\pagestyle{plain}
\chapter{Introduction}

\section{The Cosmological Context}

Neutral (atomic) hydrogen (HI) plays an important role in cosmology and structure formation processes: its distribution follows the underlying matter density field, and for most of the cosmic history it constitutes the reservoir of baryons to fuel star formation. This makes it a powerful and novel tracer of the large-scale structure (LSS) of the Universe~\citep[e.g.][]{review,Ansari_2012,Santos:2015,villa14}. While Cosmic Microwave Background (CMB) observations~\citep{WMAP,planck:2018,ACT} and galaxy redshift surveys~\citep[e.g.][]{BOSS:2016} have significantly constrained the parameters of the standard $\Lambda$CDM cosmological model, key questions are still unresolved. In particular, the fundamental nature of dark matter and dark energy is still unknown, and persistent tensions between different cosmological measurements have yet to be fully understood, including the recent tentative evidence for evolving dark energy~\citep[e.g.][]{riess:2019,Verde:2019,wong:2019,DESIde}. Mapping the large-scale distribution of HI and tracking its evolution over cosmic time offer a complementary approach to traditional galaxy surveys, by probing large volumes at high redshifts, and thus potentially providing stress tests for many cosmological models in new regimes~\citep[e.g.][]{bull:2015,villa15,obuljen_high-redshift_2018,berti_constraining_2022,Villaescusa-Navarro:2016}.

The 21 cm line, which results from the spin-flip transition within the hyperfine structure of the ground state of neutral hydrogen (see~\citet{Furlanetto:2006}), is redshifted by the cosmological expansion and can be observed at radio wavelengths. Many experiments, including interferometric arrays like CHIME~\citep{Bandura:2014gwa,CHIMEdetection}, CHORD, and HIRAX~\citep{Newburgh:2016mwi}, and single-dish instruments such as GBT~\citep{Masui2013,Wolz2022} and FAST~\citep{Hu:2019okh}, aim to detect this signal using intensity mapping (IM) techniques~\citep{Bharadwaj:2000,Battye:2004,McQuinn:2005,Chang:2007,Seo:2009fq,Battye:2013,Kovetz:2017}. Several of these efforts have already achieved detections through cross-correlation with optical galaxy surveys~\citep{Chang2010,Masui2013,Anderson2018,Wolz2022,Cunnington:2022,Paul2023}.

Radio cosmology is also a key science objective for the SKA Observatory (SKAO),\footnote{\url{https://www.skao.int/}} which will comprise two major arrays: SKA-Low in Australia and SKA-Mid in South Africa. In particular, SKA-Mid, when operated in single-dish mode~\citep[e.g.][]{Santos:2015,Bacon:2018}, will be capable of conducting 21 cm IM surveys across cosmologically relevant scales out to redshift $z \sim 3$. Currently under construction, the SKAO has a working precursor, MeerKAT, which is already contributing through its cosmological IM survey, MeerKLASS~\citep{Santos:2017}. Early results from MeerKAT data have been encouraging~\citep{Wang:2021,irfan2022}, including a first cross-correlation detection with WiggleZ galaxy data~\citep{Cunnington:2022uzo}.

Alongside several technical efforts mainly focused on foreground modeling and mitigation, refining the forecast capabilities of 21 cm IM, both as a standalone probe and in combination with other cosmological observables, is crucial (e.g.~\citep{carucci_cross-correlation_2017,berti_21cm_2023,berti24}). These forecasts are essential not only for strengthening the scientific case for 21 cm IM radio cosmology, but also in order to optimize the design and strategy of upcoming surveys.

% --- NEW: LSS-depth passage (added per supervisor request) -----------
The large-scale structure (LSS) traced by these observables is itself the end product of a well-understood but analytically intractable evolutionary history. According to the standard inflationary paradigm, the initial density perturbations that seeded structure formation were, to very high precision, a Gaussian random field, fully characterized by its two-point statistics (the power spectrum) and its early-time transfer function~\citep[e.g.][]{planck:2018}. Under gravitational instability, these perturbations grow, and while linear perturbation theory accurately describes their evolution on large scales and at early times, the growth becomes non-linear once density contrasts approach unity, first on small scales and progressively on larger ones as cosmic time advances. The result is the LSS observed today: a highly non-Gaussian density field, organized into the filamentary cosmic web of halos, voids, and filaments~\citep[e.g.][]{obuljen23,villaescusa-navarro_ingredients_2018}.

This non-Gaussianity has direct consequences for how much cosmological information can be extracted from LSS observables. For a Gaussian field, the power spectrum is a complete statistical description, and all cosmological information is contained within it. Once gravitational (and, for baryonic tracers, astrophysical) non-linearities become important, however, this is no longer true: information leaks into higher-order correlations (the bispectrum, trispectrum, and beyond), and standard perturbative techniques used to model these statistics break down at the mildly non-linear and non-linear scales that are often the most information-rich~\citep[e.g.][]{obuljen23}. Restricting analyses to the two-point function alone therefore discards a substantial fraction of the cosmological information encoded in the field, particularly on the small scales probed by upcoming high-resolution 21 cm IM and galaxy surveys.

This has motivated a shift, over the last several years, toward \emph{field-level} inference and emulation: rather than compressing a cosmological field into a small set of hand-chosen summary statistics before performing inference, one instead works directly with the field itself (or a learned representation of it), in principle retaining the full non-Gaussian information content~\citep[e.g.][]{jamieson_field-level_2023,angulo_large-scale_2022}. A representative example of this direction is the work of \citet{rouhiainen_cosmology_2024}, which explicitly frames the non-Gaussianity of the late-time, small-scale LSS as the central obstacle to analytic modeling, and addresses it by learning the field-level probability density of the matter distribution directly with normalizing flows -- explicitly analogous to how gravitational evolution transforms an initially Gaussian field into today's complex non-Gaussian one -- before turning to denoising diffusion models for volumetric, 3D super-resolution tasks. The approach taken in this thesis shares this field-level philosophy: rather than relying on summary statistics or hand-calibrated halo-occupation-style prescriptions, we aim to learn the field-level mapping between cosmological tracers directly from simulations, retaining small-scale non-Gaussian structure that summary-statistic-based or purely perturbative approaches would otherwise discard.
% --- END NEW ------------------------------------------------------

% --- NEW 2: Quantitative "Growth of Structure" / "Nonlinear LSS" depth pass ---
% Adapted from the derivation chain in Rouhiainen (2024, arXiv:2402.07694),
% Chapter 1, using YOUR macros (\Om, \sige, \hub) where they already exist.
% See "NEW BIBKEYS NEEDED" at the end of this file -- these citations are
% NOT yet in your .bib and must be added before this compiles cleanly.

\subsection{Growth of structure}

To make the qualitative picture above quantitative, it is useful to work directly with the matter overdensity field,
\begin{equation}
    \delta(\bm{r},t)\equiv\frac{\rho(\bm{r},t)-\langle\rho(\bm{r},t)\rangle}{\langle\rho(\bm{r},t)\rangle},
\end{equation}
where $\langle...\rangle$ denotes an ensemble average. On sub-horizon scales and at times well after recombination, $\delta$ obeys, in Fourier space,
\begin{equation}
    \frac{\text{d}}{\text{d}t}\left(a^2(t)\frac{\text{d}}{\text{d}t}\delta(\bm{k},t)\right)=4\pi G a^2(t)\,\rho(t)\,\delta(\bm{k},t),
\end{equation}
sourced by the Poisson equation $k^2\Phi = \tfrac{3}{2}\Om \hub^2\,\delta(\bm{k},a)/a$. It is standard to factorize the time and scale dependence of the linear solution by introducing the growth factor $D_+(a)$ and transfer function $T(k)$,
\begin{equation}
    \delta(\bm{k},a) = \frac{2k^2}{5\Om \hub^2}\,\mathcal{R}(\bm{k})\,T(k)\,D_+(a),
\end{equation}
where $\mathcal{R}(\bm{k})$ encodes the primordial (inflationary) fluctuations. The growth factor satisfies $D_+(a)\propto a$ (matter domination) at early times, and is suppressed at late times by the cosmological constant, decaying relative to $a$ as dark-energy domination sets in -- explicitly,
\begin{equation}
    D_+(a) = a \; {}_2F_1\!\left(\frac{1}{3},\,1;\,\frac{11}{6};\,-\frac{\Omega_\Lambda}{\Om}a^3\right),
\end{equation}
in terms of the hypergeometric function ${}_2F_1$~\citep[e.g.][]{Weinberg2008}.

The (linear) matter power spectrum, defined through
\begin{equation}
    (2\pi)^3 \Pk\,\delta_\text{D}^3(\bm{k}+\bm{k}') = \langle\delta(\bm{k})\delta(\bm{k}')\rangle,
\end{equation}
inherits this factorization as $P_\text{Lin}(k,a) \propto T^2(k)\,D_+^2(a)$. Crucially, \emph{for a Gaussian field, the power spectrum is a complete statistical description}: every higher-order correlator $\langle\delta(\bm{k}_1)\cdots\delta(\bm{k}_n)\rangle$ either vanishes (odd $n$) or reduces to a sum of products of two-point functions (even $n$), a consequence of Wick's theorem. The amplitude of this power spectrum is conventionally summarized by $\sige$, the RMS overdensity in spheres of radius $8\,h^{-1}\,\mathrm{Mpc}$,
\begin{equation}
    \sige^2(a) = \frac{1}{2\pi^2}\int_0^\infty \text{d}k\; k^2\, P_\text{Lin}(k,a)\, W^2(k R),\qquad R = 8\,h^{-1}\,\mathrm{Mpc},
\end{equation}
with $W$ a top-hat filter in Fourier space.

\subsection{Nonlinear large-scale structure}

The linear description above breaks down once fluctuations become of order unity. This defines a nonlinear scale $k_\text{NL}$ through the dimensionless power spectrum,
\begin{equation}
    \Delta^2(k_\text{NL}) \equiv \frac{k_\text{NL}^3}{2\pi^2}\,P_\text{Lin}(k_\text{NL}) \sim 1,
\end{equation}
which for present-day $\Lambda$CDM parameters occurs at $k_\text{NL}\sim0.2\,h/\mathrm{Mpc}$~\citep[e.g.][]{Takahashi_2012} -- squarely within the regime our field-level ML models are designed to capture. Beyond this scale, non-Gaussianity manifests directly as non-vanishing higher-order statistics, most simply the bispectrum,
\begin{equation}
    (2\pi)^3 B(k_1,k_2,k_3)\,\delta_\text{D}^3(\bm{k}_1+\bm{k}_2+\bm{k}_3) = \langle\delta(\bm{k}_1)\delta(\bm{k}_2)\delta(\bm{k}_3)\rangle,
\end{equation}
which vanishes identically for a Gaussian field but carries additional, independent cosmological information once gravitational (and, for baryonic tracers, astrophysical feedback) non-linearities are significant~\citep[e.g.][]{Sefusatti_2006}. Standard cosmological perturbation theory can, in principle, extend the analytic description into this regime order-by-order (e.g.\ at ``1-loop'' order), but the required loop integrals grow rapidly in complexity, and even at next-to-leading order the perturbative bispectrum already departs from full $N$-body results by an order of magnitude on the mildly non-linear scales relevant to this work~\citep[e.g.][]{Chudaykin_2020}. This breakdown of analytic and perturbative methods precisely on the scales that carry the most cosmological information is the central motivation, throughout this thesis, for turning to machine-learning-based, field-level approaches instead.
% --- END NEW 2 ------------------------------------------------------


\section{The Simulation Bottleneck}

Cosmological simulations are essential tools for understanding the formation and evolution of structure in the Universe. They allow us to explore the complex interplay of gravitational, hydrodynamic, and astrophysical processes shaping the distribution of matter from large down to highly non-linear galactic and sub-galactic scales. However, high-fidelity hydrodynamical simulations are computationally very expensive, limiting their applicability when the parameter space is large in inference studies. For example, the IllustrisTNG simulations~\citep{nelson2021illustristngsimulationspublicdata,Pillepich_2017} required approximately 35 million CPU-hours for the largest \TNG run and about 18 million CPU-hours for TNG100, while EAGLE~\citep{EAGLE} consumed roughly 4.5 million CPU-hours.

Such computational demands become particularly prohibitive for cosmological or astrophysical field-level inference and simulation-based inference (SBI) frameworks, which require large ensembles of conditional forward simulations to explore the parameter spaces. Recent works have shown promising applications, including the cosmological context~\citep{angulo_large-scale_2022,jeffrey_likelihood-free_2020,bairagi2025simulationsneedsimulationbasedinference,Zeghal_2025,karchev2023simsimssimulationbasedsupernovaia,Saxena_2024,tucci24}. However, their applicability remains limited by the requirement of generating sufficient high-fidelity simulations. This computational bottleneck underlines the need for fast, accurate surrogate models capable of reproducing cosmological fields at a fraction of the cost — not limited to dark matter, but crucially including baryonic observables that can be directly constrained with data.

This same bottleneck also constrains theoretical modeling of the 21 cm signal itself, which currently follows three main approaches: halo models (see~\citet{padmanabhan_hi_2023}), perturbation theory (e.g.~\citep{obuljen23}), or hydrodynamical simulations of structure formation incorporating the relevant physical processes (e.g.~\citep{villaescusa-navarro_ingredients_2018}). Each of these methods has both advantages and disadvantages. For example, modeling the IM signal requires both large volumes and high resolution to fully resolve the physics of the small mass ($\sim$ 10$^{10}$ M$_{\odot}/h$) dark matter halos that host HI, something that is difficult to achieve with full hydrodynamical simulations due to the range of scales involved. In the context of halo models, extensive work has been expended to study the mass limit for halos to retain HI, for example by looking at the damped Lyman-$\alpha$ systems statistics, both from the point of view of halo models~\citep{Padmanabhan_2016} and observationally~\citep{Dev_2023,Obuljen_2019}. In addition to these semi-analytical and perturbative frameworks, several forward-modeling approaches have been developed to predict the HI field directly from dark-matter simulations by prescribing or calibrating the $M_{\rm HI}(M)$ relation and other parameters on observations or hydrodynamical runs~\citep[e.g.][]{hitz2025fastsimulationcosmologicalneutral,li2022theoreticalmodelsatomichydrogen,spinelli_atomic_2020}. These methods, however, extend predictions to large volumes by applying calibrated prescriptions for the HI--halo connection, rather than learning the full field-level mapping from simulations. Thus, a method capable of learning the small-scale physics from state-of-the-art simulations while at the same time reaching large scales and volumes would be of great importance — this is the computational and modeling gap this thesis addresses.


\section{Machine Learning for Field-Level Cosmology}

In recent years, machine learning techniques, and in particular generative models, have emerged as powerful emulators to reproduce cosmological and astrophysical fields~\citep{perraudin2019cosmologicalnbodysimulationschallenge,Li_2021,jamieson_field-level_2023,Rigo_2025,hassan_hiflow_2022,andrianomena_emulating_2022,sharma25,bernardini2025ember2,andrianomena2024cosmologicalmultifieldemulator}. Among them, diffusion models — a class of stochastic generative models that generate data by learning to reverse a gradual noising process, iteratively denoising corrupted data samples — have gained significant attention due to their ability to produce high-fidelity, diverse outputs~\citep{sohl-dickstein_deep_2015,ho_denoising_2020,kingma_variational_2023}. These models have been successfully applied in astrophysics and cosmology, for example to emulate the baryonic processes affecting galaxy properties~\citep{Chadayammuri_2023}, or more generally to augment the resolution of simulations by using a smaller set of high-resolution ones~\citep{kodi20}.

In the specific context of 21 cm IM theoretical modeling, several recent efforts aim at predicting the brightness temperature signal at the field level, including fast generative models of HI maps using normalizing flows~\citep{hassan_hiflow_2022}, the neural approach of \citet{wadekar_hinet_2021}, and a generative adversarial framework introduced by \citet{andrianomena_emulating_2022}. Diffusion models have thus far found very limited application in this domain~\citep{ono_debiasing_2024}, and none specifically within 21 cm IM. To address this, we develop a novel, physically motivated ML pipeline (\LODI) which, in a first step, uses a U-Net to learn the mapping between dark matter and halos, and in a second step employs a custom-designed diffusion model to predict the 21 cm signal down to very small, non-linear scales. Predicting both halos and the 21 cm signal jointly is crucial for using such machinery in cross-correlation studies of the IM signal with other LSS tracers (e.g.\ lensing, galaxies), which is the ultimate aim of this line of work. This should be viewed as the first demonstration of a diffusion-based, halo-conditioned approach to 21 cm IM, complementary to existing forward-modeling methods.

Despite these advances, a major challenge remains in scaling generative models to arbitrarily large, high-fidelity cosmological volumes, especially when constrained by GPU memory and computational cost. Moreover, the evolution of cosmic structure is inherently non-local: long-range gravitational interactions couple distant regions, and baryonic feedback processes — such as those driven by supernovae or active galactic nuclei — can redistribute matter over large scales. Accurately modeling these effects requires a network to learn global dependencies rather than simple local conditional relationships, which becomes particularly challenging when no similar low-resolution reference field is available to guide the generation.

To overcome these limitations, we introduce \FOLD (``Cosmological Field Overlap Latent Diffusion''), a technique building on the variational diffusion model framework \LODI presented above and in~\cite{mishra2025largefastaccuratehi}. \FOLD enables the rapid and accurate generation of large three-dimensional cosmological fields by conditioning the denoising process on the underlying large-scale matter distribution, represented by either the stellar density or dark matter density fields. We demonstrate the capability of this approach by generating gas temperature fields conditioned on dark matter density, as well as dark matter density fields reconstructed from sparse stellar density distributions, using the IllustrisTNG--TNG300 simulations. These two applications are particularly challenging and relevant since both are intrinsically sensitive to small-scale physical processes (star formation and galactic feedback) that also affect intermediate cosmological scales. Their importance is also reflected in summary statistics such as the power spectrum: it is well known that baryonic processes, as simulated in state-of-the-art hydrodynamical simulations, produce a suppression of power at $k \sim 1\,h/\mathrm{Mpc}$ caused by feedback, and a steep increase of power at smaller scales $k \gg 1\,h/\mathrm{Mpc}$ from stars forming in the innermost parts of dark matter haloes~\citep{chisari19,schneider15,vandaalen20}. While the exact amount of suppression and increase of power varies among different simulation suites, the shape of these effects is a common prediction across most models — making the temperature and stellar fields a natural choice for testing our method's performance across both regimes.

\FOLD is trained on just $\approx 1\%$ and validated on $\approx 0.5\%$ of the simulation volume, before being applied to synthesize complete large-scale realizations, with excellent performance across 1-, 2-, and 3-point statistics. It naturally enforces periodic boundary conditions and surpasses previous overlapping patch-based approaches by eliminating edge artifacts and maintaining structural coherence, even when conditioned on highly sparse stellar input.

A related, complementary goal — currently hampered by the same computational bottlenecks — is to complement traditional summary statistics (such as two- and three-point statistics) with more informative, data-driven summaries that better capture non-Gaussian features of evolved non-linear fields (e.g.~\citep{bairagi2025patchnethierarchicalapproachneural,Borg,lanzieri2025optimalneuralsummarisationfullfield,Bayer_2025}).


\section{Thesis Outline}

This thesis presents two complementary contributions to field-level, diffusion-based emulation of cosmological and astrophysical fields. [Placeholder — adapt to match final thesis chapter numbering, e.g.:] Chapter~\ref{ch:lodi} (\LODI) describes the halo-conditioned diffusion pipeline for predicting the 21 cm signal, including the U-Net architecture, diffusion model design, training simulations, loss function, and validation against held-out dark matter maps in terms of illustrative maps and power spectra. Chapter~\ref{ch:cosmofold} (\FOLD) extends this framework to large-volume field generation via latent overlap diffusion, applying it to gas temperature and dark matter density reconstruction on IllustrisTNG--TNG300. We conclude in Chapter~\ref{ch:conclusion} with a discussion of promising avenues for future work, including applications to cross-correlation studies between 21 cm IM and other LSS tracers.






\section{The Cosmological Context}
\section{The Simulation Bottleneck}
\section{Machine Learning for Field-Level Cosmology}
\section{Thesis Outline}
% \section{List of Publications}